\documentclass[12pt,a4paper]{ctexart}
\usepackage{geometry}
\usepackage{array}
\usepackage{booktabs}
\usepackage{longtable}
\usepackage{tabularx}
\usepackage{xcolor}
\usepackage{enumitem}
\usepackage{amsmath}
\usepackage{tikz}
\usepackage{hyperref}
\usepackage{fancyhdr}
\usetikzlibrary{arrows.meta,positioning,fit,shapes.geometric}

\geometry{left=2.8cm,right=2.8cm,top=2.6cm,bottom=2.6cm}
\setlength{\parindent}{2em}
\setlength{\parskip}{0.35em}
\renewcommand{\arraystretch}{1.35}
\setlist[itemize]{leftmargin=2em,itemsep=0.2em,topsep=0.2em}
\setlist[enumerate]{leftmargin=2em,itemsep=0.2em,topsep=0.2em}

\pagestyle{fancy}
\fancyhf{}
\fancyhead[C]{\small 基于时序可信动态图谱的新能源锂电池产业链黑天鹅风险推理与预警系统}
\fancyfoot[C]{\thepage}
\renewcommand{\headrulewidth}{0.4pt}

\definecolor{mainblue}{RGB}{27,79,140}
\definecolor{lightblue}{RGB}{232,242,252}
\definecolor{lightgray}{RGB}{246,247,249}

\hypersetup{
    colorlinks=true,
    linkcolor=mainblue,
    urlcolor=mainblue
}

\begin{document}

\begin{center}
    \vspace*{0.8cm}
    {\zihao{1}\bfseries 技术可行性报告\par}
    \vspace{0.8cm}
    {\zihao{3}\bfseries 基于时序可信动态图谱的新能源锂电池产业链黑天鹅风险推理与预警系统\par}
    \vspace{0.5cm}
    {\zihao{-4} 面向新能源锂电池产业链风险研判场景的技术路线、系统架构与验证方案\par}
\end{center}

\vspace{0.6cm}

\section{要解决什么问题}

新能源锂电池产业链是一个高度全球化的多层结构：上游锂、镍、钴、石墨等矿产分布在智利、澳大利亚、印尼、刚果金等国家，中游的正极材料、电芯制造集中在中国、韩国、日本，下游整车和储能市场遍布全球。这条链条上的任何一环出现政策突变、矿区停产、航运阻塞或贸易制裁，都会沿产业链快速传导，产生连锁反应。

现有手段处理这类风险时，普遍存在三个痛点：

\textbf{第一，信息过载且真假难辨。} 每天有大量关于锂矿罢工、镍出口禁令、航运中断的新闻涌入。但问题不是信息太少，而是太多且太杂。同一件事可能同时出现在政府官网、路透社报道、推特个人账号和微信公众号上——可信度完全不在一个量级。2023 年曾经有一条关于"刚果金暂停钴出口"的推文被多家资讯平台自动抓取推送，实际上只是一家地方小报引用了匿名消息，三天后才被官方澄清。如果把这类信息直接当风险信号推送给决策者，造成的误判比漏报更严重。我们需要的不是更多的信息，而是判断信息可信度的能力。

\textbf{第二，静态知识库跟不上风险演化。} 风险事件天然有一个生命周期：先是零散传闻，然后被某些渠道确认，接着可能升级（事态扩大），也可能缓解（谈判取得进展），最后要么解决要么证伪。整个过程少则几天，多则几周。普通知识图谱构建好以后就固定了，不会因为"刚果金钴政策发生了新变化"而去更新之前的记录。RAG 系统稍微好一点——能检索最新文本——但它同样不知道"这条关于刚果金钴的新闻说的是新事件还是旧闻的最新进展"。系统没有时间感知能力，就无法判断一个风险当前是正在发生、已经过去还是正在升级。

\textbf{第三，新闻文本本身不包含产业链影响路径。} 一篇路透社关于"智利锂矿工人罢工"的报道会写矿区位置、罢工原因、工会诉求、预计持续时间。但它不会自动告诉你：智利的锂占全球供应量约 30\%，其中大部分是盐湖卤水提锂，如果罢工持续超过两周，碳酸锂库存会开始紧张，进而影响正极材料厂商的采购成本，最终传导到电芯和电池包的出厂价格。这种传导关系不在新闻正文里，它藏在产业链的结构知识中。靠文本检索或大模型阅读理解都无法可靠地还原这条路径——文本里根本没有写。

我们的判断很直接：解决这三个问题，不能靠往大模型里塞更多文本。大模型擅长的是文本理解和生成，但它没有产业链结构知识，不了解我们的来源信誉体系，也计算不了图上的传播路径。我们的路线是：把\textbf{非结构化信息转化为结构化事件}，用\textbf{产业链图谱建模传导关系}，用\textbf{置信度评估过滤信息噪声}，用\textbf{图搜索计算传播路径}，最后让大模型在证据链的约束下生成报告。大模型在这个架构里是\textbf{受控的报告生成器}，不是风险判决器。

\section{三个核心难点}

在确定技术路线之前，我们把问题拆解为三个必须攻克的难点。

\subsection{难点一：多源异构信息如何统一成可计算的结构}

风险信息来源极其混杂：政府公告格式规范但语言严肃，行业新闻篇幅较长且夹带主观判断，社交媒体碎片化且真假掺杂，市场价格数据是纯数字序列。这些信息的时间粒度、可信等级、表达方式完全不同。如果每个来源各自为政，后续的推理和评分根本无法统一处理。

\textbf{解决思路：} 不是把所有来源塞进同一个模型去"理解"，而是设计一个统一的事件抽象层。无论输入来自 RSS 新闻、API 政策公告还是演示案例数据，系统都将其转化为标准化的风险事件对象。这个对象有固定的字段结构：事件类型、时间、地点、涉及的实体和材料、来源、原始文本、置信度等。后续所有模块——图谱映射、可信度计算、路径搜索、评分——都只操作这个对象，不再回头处理原始文本。

这样做的一个实际好处是：各个模块之间通过事件对象解耦。事件抽取模块可以独立替换升级（比如从规则+LLM换成微调模型），不影响下游模块。另一个好处是调试方便——任何一个中间环节出了问题，只需要检查事件对象的字段值，不需要重新跑一遍整个 NLP pipeline。

\subsection{难点二：产业链风险传播无法通过文本匹配获得}

文本检索能回答"有哪些关于锂的新闻"，但回答不了"锂矿罢工除了影响锂本身，还会影响哪些材料和环节"。后者需要产业链知识图谱——系统必须知道锂$\to$碳酸锂$\to$正极材料$\to$电芯$\to$电池包这条路径，以及其中每个节点的上下游关系。

\textbf{解决思路：} 自建面向新能源锂电池产业链的知识图谱。我们定义了五类节点——资源国（智利、澳大利亚、印尼等）、材料（锂、碳酸锂、镍、钴、石墨等）、产业环节（正极材料、电芯、电池包、整车、储能）、企业（宁德时代、比亚迪等）、物流通道（红海、马六甲等）——以及四类关系边（supplies、depends\_on、processes、transports）。种子图谱规模不大，但覆盖了从矿产到终端产品的核心传导路径。事件进入系统后先通过实体链接挂接到图谱节点，然后从节点出发沿关系边做有限深度的图搜索，找出完整的传播链路。

选择自建图谱而不是复用已有的通用知识图谱（如 Wikidata），原因是通用图谱的产业链关系粒度不适合风险推理。Wikidata 知道"锂是一种化学元素"，但不知道"锂是生产碳酸锂的原料，碳酸锂是正极材料的关键成分"。这种行业特定的上下游关系，只有面向场景自建的图谱才能提供。

\section{系统架构与整体流程}

基于上述三个难点，系统的整体架构如下：

\begin{center}
\begin{tikzpicture}[
    node distance=0.75cm,
    module/.style={rectangle, rounded corners=3pt, draw=mainblue, fill=lightblue, align=center, minimum width=10.8cm, minimum height=0.8cm, font=\small},
    arrow/.style={-{Latex[length=2.2mm]}, thick, draw=mainblue}
]
\node[module] (a) {多源信息输入：新闻、政策公告、行业媒体、市场指标、演示案例};
\node[module, below=of a] (b) {风险事件结构化：类型、时间、地点、实体、材料、证据、严重度};
\node[module, below=of b] (c) {产业链知识图谱映射：国家、资源、材料、企业、物流通道};
\node[module, below=of c] (d) {时序可信动态图谱推理：时间状态、置信度、冲突识别、路径搜索};
\node[module, below=of d] (e) {风险传播与评分：严重度、节点关键性、来源置信度、传播范围};
\node[module, below=of e] (f) {Dashboard 态势展示：地图点位、传播弧线、事件状态、市场指标};
\node[module, below=of f] (g) {AI 研判报告生成：基于结构化证据链生成穿透式风险报告};
\draw[arrow] (a) -- (b);
\draw[arrow] (b) -- (c);
\draw[arrow] (c) -- (d);
\draw[arrow] (d) -- (e);
\draw[arrow] (e) -- (f);
\draw[arrow] (f) -- (g);
\end{tikzpicture}
\end{center}

整个链条的设计原则是：\textbf{每一步都在增加信息的结构化程度，每一步都在压缩后续模块的判断空间。} 事件结构化把自由文本变成定长字段，图谱映射把文本实体固定到已知节点上，可信度计算给每个事件打了置信度标签，路径搜索限定传播范围，评分模型给出量化分数——到大模型这一步，输入已经是一份结构清晰的证据包，而不是一篇需要"阅读理解"的新闻。

这七个层级的具体职责如下：

\textbf{第一层 多源信息输入：} 系统不依赖单一类型的信息源。新闻类源（通过 RSS 或 API 接入路透社、行业媒体等）提供最新事件信息。政策公告类源（政府官网、国际组织发布页）提供高可信度的官方信息。市场数据类源（锂矿股价格、碳酸锂现货价、VIX 指数等）提供市场对风险的实时反应信号。演示案例数据是预置的 30--50 条典型场景事件，确保在网络不可用或源不可达时系统仍能完整演示。所有信息源进入系统时都会附带来源元数据（来源名、类型、信誉等级），供后续可信度评估使用。

\textbf{第二层 风险事件结构化：} 这是把"新闻"变成"数据"的关键步骤。原始自然语言文本经过规则+LLM 混合管线处理，输出标准化的风险事件对象。对象包含结构化字段（类型、时间、地点、实体、严重度等），使得后续模块不需要理解自然语言，只需要读取字段值。这一层也是整个系统唯一需要 NLP 能力的地方，后续所有模块都是结构化数据处理。

\textbf{第三层 产业链知识图谱映射：} 事件对象中的实体（如"印尼""镍""出口限制"）需要挂接到产业链图谱的已知节点上。映射过程使用四级回退策略（从精确别名匹配到图谱关联推断），确保即使文本表述不规范的新闻也能找到对应的图谱位置。映射完成后，事件就获得了在图谱中的"坐标"——它影响哪个节点，以及该节点在图谱中的上下游关系是什么。

\textbf{第四层 时序可信动态图谱推理：} 这一层给事件增加了两个关键维度：时间和可信度。时间维度通过事件状态机（unverified$\to$watchlist$\to$confirmed/conflicting$\to$resolved）追踪事件的生命周期，使系统能区分"三天前发生的事件"和"今天有新进展的事件"。可信度维度通过置信度计算公式，结合来源信誉、多源一致性和证据一致性为每个事件打分，使系统能区分"确认的事实"和"未经证实的传闻"。

\textbf{第五层 风险传播与评分：} 这一层回答两个问题：这个风险会影响到产业链的哪些环节（传播路径），以及影响有多大（风险分数）。传播路径通过图搜索算法从风险源节点出发遍历整个产业链图谱找到受影响节点链。风险分数通过五因子加权模型给出 0--100 的量化评分，每个因子都有明确的业务含义和计算方式。

\textbf{第六层 Dashboard 态势展示：} 后端推理的所有结果汇总到统一的 Dashboard 数据接口。前端根据数据渲染全球地图点位、传播弧线、事件列表、市场指标面板和 AI 报告面板。关键设计理念是"把中间结果也展示出来"——不仅仅是最终的风险分数，还包括置信度状态、数据缺口、推断标记等，让用户能审查系统的判断依据。

\textbf{第七层 AI 研判报告生成：} 所有结构化数据（事件、评分、路径、证据）组装成 Prompt，输入大语言模型生成自然语言报告。大模型的自由度被严格限制——它只能基于输入的结构化数据组织语言，不能自由发挥或补充信息。报告输出后经过格式校验，不合格则降级为模板填充的 fallback 报告。

\section{核心数据对象设计}

\subsection{风险事件对象}

风险事件对象是系统推理的基本单元。系统将新闻、公告和演示案例统一转化为标准化事件对象，使后续可信度计算、实体映射、图谱推理和报告生成具有统一输入。

事件对象采用 JSON 格式定义，包含六个字段类别，每个类别回答一类问题：

\begin{table}[htbp]
\centering
\caption{风险事件对象字段设计}
\begin{tabularx}{\textwidth}{p{2.0cm}p{2.8cm}p{3.8cm}X}
\toprule
字段类别 & 代表字段 & 技术作用 & 值示例 \\
\midrule
事实字段 & event\_name、event\_type、published\_at、location & 描述事件基本事实——发生了什么、何时发生、在哪里发生。 & "智利锂矿工人罢工"、production\_disruption、2026-06-15、智利 \\
产业链字段 & materials、affected\_nodes、matched\_nodes & 将事件与材料、企业、国家和产业环节连接，供图谱映射使用。 & ["锂"]、["智利"], ["锂\_node"] \\
证据字段 & source、evidence、confidence & 支撑可信度计算、证据引用和报告追溯。 & Reuters、原文引用URL、0.85 \\
时序字段 & event\_time、valid\_from、valid\_to、lifecycle & 判断事件是否仍然有效，以及风险处于什么阶段。 & 2026-06-15、2026-06-15、2026-07-15、active \\
风险字段 & severity、risk\_score、risk\_level & 支撑风险排序、等级划分和预警展示。 & 4、73、high \\
解释字段 & propagation\_paths、explanation & 传播路径解释、Dashboard 展示和人工复核。 & ["锂→碳酸锂→正极材料→电芯"]、"智利锂矿罢工预计影响碳酸锂供应，沿产业链传导至正极材料和电芯环节。" \\
\bottomrule
\end{tabularx}
\end{table}

一个完整的事件对象在系统中的实际表示如下（简化）：

\begin{verbatim}
{
  "event_id": "EVT-20260615-001",
  "event_name": "智利锂矿工人罢工",
  "event_type": "production_disruption",
  "published_at": "2026-06-15T09:30:00Z",
  "location": { "name": "智利", "lat": -23.65, "lng": -70.40 },
  "materials": ["锂"],
  "source": { "name": "Reuters", "type": "authoritative_media", "reputation": 0.88 },
  "confidence": 0.85,
  "lifecycle": "active",
  "severity": 4,
  "risk_score": 73,
  "risk_level": "high",
  "propagation_paths": [
    "锂→碳酸锂→正极材料→电芯→电池包"
  ],
  "explanation": "智利盐湖提锂产量占全球约30%，罢工若持续超过两周..."
}
\end{verbatim}

该设计使系统能够回答六类问题：事件是什么（事实字段）、来自哪里（证据字段）、是否可信（置信度）、影响什么（产业链字段）、如何传播（传播路径）、风险多高（评分字段）。相比普通新闻摘要，事件对象是结构化数据——可排序、可过滤、可做算术运算、可沿图谱路径搜索。

\subsection{产业链图谱对象}

产业链图谱由节点和边组成。节点表示国家、资源、材料、企业、物流通道、产业环节和需求场景；边表示供应、依赖、加工、运输、影响和替代关系。

\begin{table}[htbp]
\centering
\caption{图谱节点与边设计}
\begin{tabularx}{\textwidth}{p{2.6cm}p{4.6cm}X}
\toprule
图谱元素 & 示例 & 作用 \\
\midrule
资源国节点 & 智利、澳大利亚、阿根廷、印尼、刚果金 & 表示关键矿产资源来源。 \\
材料节点 & 锂、碳酸锂、氢氧化锂、镍、钴、石墨 & 表示上游和中游材料。 \\
产业环节节点 & 正极材料、电芯、电池包、整车、储能 & 表示上下游传导结构。 \\
企业节点 & 电池企业、整车企业、材料企业 & 表示潜在受影响主体。 \\
物流节点 & 红海、苏伊士运河、马六甲海峡 & 表示跨境运输风险通道。 \\
关系边 & supplies、depends\_on、processes、transports & 表示传播方向和依赖关系。 \\
\bottomrule
\end{tabularx}
\end{table}

图谱设计的关键不是节点数量越多越好，而是要覆盖风险传播所需的最小可用产业链结构。初赛阶段的种子图谱可支撑典型案例演示，后续可继续扩展企业、矿山、供应商和贸易关系。

\section{关键方法的设计与实现}

\subsection{事件结构化：规则兜底，大模型增强}

事件抽取面临一个实际矛盾：大模型理解能力强但输出不稳定——同样的新闻每次调用返回的 JSON 字段名可能不一样，甚至有时候会漏掉关键字段。规则系统稳定但覆盖范围有限——写死的正则和词典只能处理已知的风险类型和实体名称。

我们的做法是\textbf{规则做前过滤和后校验，大模型做中间补全}，具体分四步：

\textbf{第一步，候选实体识别（规则）。} 原始文本进来后，先用正则和词典匹配材料词（锂、碳酸锂、镍、钴等 60+ 个词条）、国家/地区词（智利、印尼、刚果金等 30+ 个）、企业词（头部电池和车企 20+ 个）、风险类型词（罢工、禁令、制裁、中断等 15+ 类）。这一步不追求召回率，但要求精度——宁可漏掉一些模糊表达，也不能把无关词误判为候选实体。

\textbf{第二步，事件类型分类（规则+关键词）。} 根据匹配到的风险类型词和文本中的触发词（"宣布""暂停""限制""警告"等），判断事件属于政策变化、生产中断、物流中断、价格波动、自然灾害、地缘冲突中的哪一类。如果文本中同时出现多种风险信号，取置信度最高的一类。

\textbf{第三步，字段补全（LLM）。} 把原始文本和第一步得到的候选实体列表一起塞给大模型，让它在候选列表的约束下补全事件的其他字段：精确的事件时间、影响描述、严重度评估、涉及的上下游材料和环节。候选实体列表的作用是给模型划定边界——它只能从列表里选实体，不能自己发明。这一步是唯一依赖大模型的地方。

\textbf{第四步，Schema 校验（规则）。} 大模型返回的结果经过 schema 校验，检查每个字段的格式和合法性：时间字段是否是合法日期，严重度是否在 1--5 的范围内，实体是否在候选列表或图谱节点中。校验不通过的字段触发第二次 LLM 补全，只传回缺失或不合法的字段让模型重新生成，而不是从头再来。

\begin{center}
\begin{tikzpicture}[
    node distance=0.55cm,
    snode/.style={rectangle, draw=mainblue, fill=lightgray, align=center, minimum width=3.0cm, minimum height=0.8cm, font=\small},
    arrow/.style={-{Latex[length=2mm]}, thick, draw=mainblue}
]
\node[snode] (s1) {原始文本};
\node[snode, right=of s1] (s2) {候选实体识别};
\node[snode, right=of s2] (s3) {事件类型分类};
\node[snode, below=of s3] (s4) {字段补全（LLM）};
\node[snode, left=of s4] (s5) {Schema 校验};
\node[snode, left=of s5] (s6) {图谱节点映射};
\node[snode, below=of s6, minimum width=9.8cm] (s7) {标准风险事件对象};
\draw[arrow] (s1) -- (s2);
\draw[arrow] (s2) -- (s3);
\draw[arrow] (s3) -- (s4);
\draw[arrow] (s4) -- (s5);
\draw[arrow] (s5) -- (s6);
\draw[arrow] (s6) -- (s7);
\draw[arrow] (s4) -- (s5);
\end{tikzpicture}
\end{center}

这样设计的好处是：规则保证了输出格式的稳定性——即使大模型调用超时或返回乱码，规则模块仍然能输出一个结构完整的事件对象（虽然可能缺严重度等字段）。大模型只在需要"理解"的地方介入——补全描述、归纳影响、评估严重度。整个管线不是"先大模型再人工校验"，而是"规则确定框架，大模型填空，规则再检查填空质量"。

\subsection{实体链接：四级逐级回退策略}

事件抽取完成后，面临的现实问题是：新闻文本中的实体表述方式千奇百怪——"DRC"和"刚果金"指向同一个国家，"Red Sea"和"红海"指向同一个物流通道，有时候文本里只写"cathode"不写"正极材料"。如果实体链接在这一步失败，后续所有图谱推理都会落空。

我们维护了一个别名表来解决这个问题。别名表是手工整理的，包含三类映射：中英文名称对照（nickel$\leftrightarrow$镍、cobalt$\leftrightarrow$钴）、全称缩写映射（DRC$\leftrightarrow$刚果金）、行业术语映射（cathode$\leftrightarrow$正极材料、三元$\leftrightarrow$镍钴锰酸锂）。当前版本约 150 条映射条目，覆盖了产业链核心实体最常见的表述变体。

但别名表不可能覆盖所有情况——新闻里可能会写"南美主要锂生产国"而不是"智利"。因此我们设计了四级回退策略：

\begin{table}[htbp]
\centering
\caption{实体链接策略}
\begin{tabularx}{\textwidth}{p{1.8cm}p{3.0cm}p{4.0cm}X}
\toprule
层级 & 匹配方法 & 示例 & 适用场景 \\
\midrule
第一级 & 精确别名匹配 & Red Sea $\to$ 红海物流通道 & 文本中出现实体标准名或别名表中登记的变体。 \\
第二级 & 材料/行业词匹配 & nickel $\to$ 镍 & 文本中出现材料名称或英文行业术语。 \\
第三级 & 行业上下文推断 & cathode $\to$ 正极材料 & 文本有明确行业语境但缺乏标准节点名。 \\
第四级 & 图谱关联回退 & lithium $\to$ 智利、澳大利亚 & 文本中有材料但缺精确位置时做可解释推断。 \\
\bottomrule
\end{tabularx}
\end{table}

第四级回退会明确标记为"inferred"而非"exact"。这个标记会沿着管线传递——在传播推理中，基于推断链接的路径会被降低权重；在报告生成中，推断点位会标注"数据来源：图谱推断"而不是直接写成确定事实。这样做的意图是：宁可让用户知道系统不确定，也不能把猜测包装成事实。

\subsection{可信度评估：把"该信谁"变成可计算的指标}

可信度评估是系统区别于普通新闻聚合的关键模块。我们定义事件置信度由四个因子加权合成：

\[
C_{\text{event}} = 0.40 \times C_{\text{source}} + 0.25 \times C_{\text{multi}} + 0.25 \times C_{\text{evidence}} + 0.10 \times C_{\text{entity}}
\]

各因子的设计逻辑如下：

\textbf{\(C_{\text{source}}\)（来源信誉度）：} 我们对信息来源做了分级。政府官网、国际组织公告和法院文件属于第一梯队，信誉度取值 0.90--0.98。路透社、行业研究机构等权威媒体属于第二梯队，取值 0.75--0.90。一般新闻网站和财经媒体属于第三梯队，取值 0.55--0.75。社交媒体（推特、Reddit、微信公众号）属于第四梯队，取值 0.30--0.55。取值是区间而非固定值——同一来源内，报道质量有差异时会浮动。

\textbf{\(C_{\text{multi}}\)（多源一致性）：} 如果一个事件被两个以上独立来源报道，且报道的核心事实（时间、地点、主体、影响）基本一致，该项得分会显著高于单来源事件。独立来源的定义是来源之间没有明显的内容转载关系——比如路透社和彭博社各自派记者采访的报道算独立来源，但 A 媒体全文转载 B 媒体的稿件不算。

\textbf{\(C_{\text{evidence}}\)（证据一致性）：} 检查不同来源在关键事实维度上是否相互印证。比如一条关于"印尼镍出口限制"的消息，如果政府官网说"正在研究"，行业媒体说"已确定实施"，社交媒体说"立即生效"，三者之间存在明显矛盾，该项会扣分。

\textbf{\(C_{\text{entity}}\)（实体匹配程度）：} 事件中的实体能否精确匹配到图谱节点。精确别名匹配（第一级）得分最高，图谱关联回退（第四级）得分最低。

基于合成置信度，事件进入以下状态机：

\begin{center}
\begin{tikzpicture}[
    node distance=0.8cm,
    state/.style={rectangle, rounded corners=3pt, draw=mainblue, fill=lightblue, align=center, minimum width=3.0cm, minimum height=0.75cm, font=\small},
    arrow/.style={-{Latex[length=2mm]}, thick, draw=mainblue}
]
\node[state] (a) {单源\\unverified};
\node[state, right=of a] (b) {高可信单源\\watchlist};
\node[state, right=of b] (c) {多源一致\\confirmed};
\node[state, below=of b] (d) {多源冲突\\conflicting};
\node[state, below=of c] (e) {已解决或已失效\\resolved / invalidated};
\draw[arrow] (a) -- (b);
\draw[arrow] (b) -- (c);
\draw[arrow] (a) -- (d);
\draw[arrow] (b) -- (d);
\draw[arrow] (c) -- (e);
\draw[arrow] (d) -- (e);
\end{tikzpicture}
\end{center}

状态转换遵循几条规则：只有被至少两个独立来源交叉印证的事件才能进入 confirmed；被高信誉单源报道但未经交叉印证的事件归入 watchlist；多来源信息明显矛盾的事件归入 conflicting（此时不会参与风险评分，直到矛盾被解决）；被官方后续公告证伪或已明确解决的事件进入 resolved/invalidated。

\textbf{需要说明的是：} 当前权重值是原型阶段根据常识和行业经验设定的初始值，不是通过数据驱动方式优化的最优参数。我们承认这四个因子的权重合理性还有待验证。后续计划通过以下方式系统化调优：构建 200--300 条人工标注的验证集（每个事件标注真实置信度），然后用网格搜索和交叉验证找出最优权重组合。但即便在当前初始状态下，这个置信度框架已经比"所有来源等权处理"要好——至少它给了系统一个区分信息来源质量的机制。

\subsection{图谱传播路径推理：从"影响什么"到"如何传导"}

这是系统另一个关键模块。风险事件的影响不是孤立的——印尼限制镍出口，影响的不仅是镍这个矿种，而是沿着"镍$\to$硫酸镍$\to$三元正极材料$\to$电芯"的路径传导下去。新闻不会写出这条路径，但产业链图谱知道。

传播推理的核心是一个带权有向图搜索过程，具体分四步：

\textbf{第一步，确定风险源节点。} 根据事件类型和实体链接结果，找到图谱中对应的起始节点。一个事件可能对应多个源节点——比如"印尼限制镍出口，同时影响钴供应"，源节点就是"印尼 + 镍"和"印尼 + 钴"。源节点还附带事件严重度作为初始权重。

\textbf{第二步，沿图搜索候选影响路径。} 从源节点出发，沿 supplies（供应）、depends\_on（依赖）、processes（加工）、transports（运输）四类关系边进行搜索。搜索算法同时支持 BFS（用于快速获取广泛影响范围）和带权重的 DFS（用于评估特定路径的深度影响）。默认最大搜索深度为 5 跳——超过 5 跳的传播关系在实际产业链中已经非常微弱，纳入推理反而引入噪声。

\textbf{第三步，路径评分与排序。} 找到的每条路径按照以下因素综合排序：节点在图谱中的拓扑关键性（我们用 PageRank 和介数中心性两个指标衡量——PageRank 高的节点被大量上下游节点依赖，介数中心性高的节点是产业链中的关键枢纽）、路径长度（越短的路径影响越直接）、事件类型与节点类型的匹配度（"罢工"匹配"生产中断"类节点比匹配"物流中断"类节点得分更高）、置信度沿路径的衰减（经过的边越多，初始置信度打折扣越多，每条边衰减系数约 0.85）。

\textbf{第四步，聚合传播结果。} 所有搜索到的路径汇聚成一张风险传播子图。子图中的每个节点得到一个综合影响分数，由经过该节点的所有路径的评分之和决定。评分最高的前 N 个节点作为重点影响目标输出到下游模块。

传播推理的输出不仅仅是路径列表——还包括一张风险传播子图。这张子图会被 Dashboard 的地图组件用来画传播弧线，也会被 LLM 报告生成模块用来描述"风险如何传导"。

典型传播路径示例：

\begin{table}[htbp]
\centering
\caption{典型风险传播路径}
\begin{tabularx}{\textwidth}{p{2.6cm}p{3.0cm}X}
\toprule
事件类型 & 风险源节点 & 传播路径 \\
\midrule
矿区罢工 & 智利 / 锂 & 锂$\to$碳酸锂$\to$正极材料$\to$电芯$\to$电池包 \\
出口管制 & 印尼 / 镍 & 镍$\to$硫酸镍$\to$三元正极$\to$电芯 \\
航运中断 & 红海 & 红海$\to$欧亚航运$\to$材料交付$\to$电池生产 \\
环保监管 & 石墨 / 负极材料 & 石墨$\to$负极材料$\to$电芯$\to$储能系统 \\
\bottomrule
\end{tabularx}
\end{table}

这个方法把风险研判从"文本相似度匹配"提升到了"产业链结构推理"的层面。大模型生成报告时使用这些路径作为结构化输入，因此它能解释"为什么这个事件会影响某家企业"，而不是泛泛地说"存在供应链风险"。

\subsection{可解释风险评分：每个分数都有来源}

风险评分解决的是"这个风险有多严重"的问题。我们选择线性加权形式而非神经网络或决策树，原因是可解释性优先——评委和用户需要理解为什么一个事件得了 73 分而不是 68 分。

评分模型由五个因子加权合成：

\[
\begin{aligned}
\text{RiskScore} = &\; 0.35 \times \text{Severity} + 0.25 \times \text{Criticality} \\
& + 0.20 \times \text{Confidence} + 0.15 \times \text{Propagation} + 0.05 \times \text{Freshness}
\end{aligned}
\]

\begin{table}[htbp]
\centering
\caption{评分因子说明}
\begin{tabularx}{\textwidth}{p{2.8cm}p{1.2cm}p{4.3cm}X}
\toprule
因子 & 权重 & 具体含义 & 为什么不直接用大模型打分 \\
\midrule
事件严重度 & 0.35 & 事件本身的冲击强度。由事件结构化模块评估，分为 5 级：轻微(1)、关注(2)、重要(3)、严重(4)、危机(5)。 & 大模型容易把情绪化表述（"毁灭性打击""前所未有"）当成严重度依据，而且不同模型对严重度的判断标准不一致。 \\
节点关键性 & 0.25 & 受影响节点在图谱中的结构地位。综合 PageRank 和介数中心性归一化到 0--1。 & 需要全图拓扑计算，大模型看不到整个产业链图结构。 \\
来源置信度 & 0.20 & 证据可靠性。直接从可信度评估模块获取。 & 大模型不了解我们构建的来源信誉体系——它不知道某条新闻来自哪个层级的信源。 \\
传播影响范围 & 0.15 & 风险沿产业链扩散的广度。根据传播推理子图中受影响的不同环节数量计算。 & 需要图搜索确定扩散范围，大模型无法从单条新闻推断产业链覆盖度。 \\
时间新鲜度 & 0.05 & 事件发生距今的时效性。24 小时内满分，之后按天线性衰减，30 天后归零。 & 权重故意设得低——避免系统过度放大刚发生的小事件。 \\
\bottomrule
\end{tabularx}
\end{table}

风险等级划分：0--30 低风险（绿色），31--60 中风险（黄色），61--80 高风险（橙色），81--100 紧急风险（红色）。

这个评分模型的核心优势不在于权重是否最优（显然需要后续调优），而在于\textbf{每个分数都能拆解出具体贡献来源}。Dashboard 上展示风险分数时，会同时展示分项贡献——"73 分 = 严重度 25 分 + 节点关键性 18 分 + 来源置信度 14 分 + 传播范围 12 分 + 新鲜度 4 分"。评审或专家复核时，如果觉得分数不合理，可以直接看到哪个因子贡献过高或过低，而不是面对一个不知道从哪来的 73 分。

\subsection{证据链约束的大模型报告生成}

大模型在这个系统里的定位经过了几轮调整才最终确定。最初我们尝试过让大模型直接读新闻出报告——结果发现它经常在报告里补充事实原文中没有的信息（幻觉），而且不同批次生成的报告风格差异很大。后来我们把定位收窄为\textbf{不是风险分析师，而是报告撰写员}。

它不负责判断风险大小（评分模型已经算好了），不负责判断事件是否可信（置信度评估已经做完了），不负责找传播路径（图搜索已经算出来了）。它的任务只有一个：把结构化输入组织成一份可读的中文风险研判报告。

\begin{center}
\begin{tikzpicture}[
    node distance=0.55cm,
    snode/.style={rectangle, draw=mainblue, fill=lightgray, align=center, minimum width=2.8cm, minimum height=0.8cm, font=\small},
    arrow/.style={-{Latex[length=2mm]}, thick, draw=mainblue}
]
\node[snode] (a) {结构化事件};
\node[snode, right=of a] (b) {风险评分};
\node[snode, right=of b] (c) {传播路径};
\node[snode, below=of c] (d) {证据来源};
\node[snode, left=of d] (e) {Prompt 组装};
\node[snode, left=of e] (f) {LLM 报告生成};
\node[snode, below=of f, minimum width=8.5cm] (g) {Fallback 校验与报告输出};
\draw[arrow] (a) -- (b);
\draw[arrow] (b) -- (c);
\draw[arrow] (c) -- (d);
\draw[arrow] (d) -- (e);
\draw[arrow] (e) -- (f);
\draw[arrow] (f) -- (g);
\end{tikzpicture}
\end{center}

输入给大模型的内容是一个组装好的结构化 Prompt，包含以下部分：

\begin{itemize}
    \item \textbf{系统指令：} 定义模型角色（"你是一名产业链风险分析助理，你的任务是基于系统提供的结构化数据撰写风险研判报告，不得添加外部信息"）和输出规则。
    \item \textbf{约束条件列表：} 明确定义模型不能做的事——不得补充外部事实、不得修改系统提供的风险分数、不得将置信度低于 confirmed 的事件写成确定事实、不得跳过证据链直接给出结论。
    \item \textbf{结构化事件列表：} 每条事件包含类型、时间、地点、严重度、置信度状态。格式如："事件1: 智利锂矿罢工 | 生产中断 | 2026-06-15 | 严重度4 | confirmed"。
    \item \textbf{风险评分分解：} 每个事件的总分和各分项得分。格式如："智利锂矿罢工: 总分73=严重度28+关键性18+置信度14+传播12+新鲜度1"。
    \item \textbf{传播路径子图：} 风险源节点到受影响节点的路径链。格式如："智利/锂 $\to$ 碳酸锂 $\to$ 正极材料 $\to$ 电芯 $\to$ 电池包"。
    \item \textbf{证据来源清单：} 每条事件引用的信息来源及其信誉等级。格式如："路透社(0.88) + 智利矿业部官网(0.95) = 多源confirmed"。
    \item \textbf{数据缺口说明：} 明确告知模型哪些数据不可用（如"锂现货价格数据源当前不可用"），模型在报告中应如实反映而非回避。
\end{itemize}

实际组装后的 Prompt 结构大致如下（简化）：

\begin{verbatim}
【系统指令】
你是一名产业链风险分析助理。你的任务是基于结构化数据撰写风险研判报告。
约束：1.不得补充外部事实。2.不得修改风险分数。3.未confirmed事件不得写成确定事实。
4.每条结论必须引用证据来源。5.数据缺口必须如实说明。

【结构化事件】
- 事件1: 智利锂矿罢工 | 生产中断 | 2026-06-15 | 严重度4 | confirmed
- 事件2: 印尼镍出口限制 | 政策变化 | 2026-06-20 | 严重度3 | watchlist

【风险评分】
- 智利锂矿罢工: 73分 = 严重度28+关键性18+置信度14+传播12+新鲜度1
- 印尼镍出口限制: 52分 = 严重度21+关键性15+置信度8+传播6+新鲜度2

【传播路径】
- 智利/锂 → 碳酸锂 → 正极材料 → 电芯 → 电池包
- 印尼/镍 → 硫酸镍 → 三元正极 → 电芯

【证据来源】
- 智利事件: 路透社(0.88) + 智利矿业部官网(0.95) = 多源confirmed
- 印尼事件: 行业媒体(0.72) = 单源watchlist

【数据缺口】
- 锂现货价格数据源: 当前不可用
- 印尼镍矿实时产量: 无公开数据
\end{verbatim}

输出格式固定为 8 个部分：事件概述、影响链路、关键节点、评分分解、证据清单、应对建议、数据缺口、不确定性说明。每部分都有明确的字数限制和内容规范，不允许自由发挥。

如果大模型调用失败（超时、断连）或输出格式校验不通过（缺少必须的章节），系统自动降级为 fallback 报告。Fallback 报告不依赖大模型——它基于模板直接填充结构化数据，生成纯文本格式的风险通报。这种设计保证在任何演示场景下，即使大模型 API 完全不可用，系统仍然能输出一份包含事件信息、评分和路径的可用报告。

\section{Dashboard：不只是可视化，更是推理结果的验证界面}

Dashboard 不仅仅是一个漂亮的前端页面。它的真实作用是\textbf{把后端推理的每个中间结果都暴露给用户审查}。

后端通过 \texttt{/api/lithium/dashboard} 接口输出 9 个数据块：

\begin{table}[htbp]
\centering
\caption{Dashboard 数据接口}
\begin{tabularx}{\textwidth}{p{2.5cm}p{5.6cm}X}
\toprule
数据块 & 内容 & 做什么用 \\
\midrule
summary & 事件总数、最高风险、状态统计 & 全局态势概览 \\
events & 结构化事件列表 & 每一条风险事件的可追溯详情 \\
map & 地图点位 + 传播弧线 & 风险位置和传导方向的可视化 \\
markets & ETF、锂矿股、VIX、能源价格 & 市场对风险的实时反应 \\
report & 缓存报告 / LLM 报告 / fallback 报告 & 三种报告级别保证展示不中断 \\
data\_quality & 数据缺口 + 来源状态 & 明确告知哪些信息不可得 \\
llm & 模型提供方 + 调用状态 & 展示技术栈透明度 \\
registry & 来源信誉注册表 & 每个信息来源的可信等级 \\
top\_events & 风险排行榜 & 按评分排序的重点关注事件 \\
\bottomrule
\end{tabularx}
\end{table}

其中 data\_quality 模块很关键——它专门用来展示"系统不知道什么"。外部数据不可用、推断点位标记、低置信度事件等都会在这里列出来。这样做的目的是：\textbf{把不确定性显式呈现，而不是藏起来}。

Dashboard 的实际数据流如下：后端将所有推理结果汇总为一个 JSON 响应，前端按数据块分别渲染。一个简化的接口响应示例如下：

\begin{verbatim}
{
  "summary": { "total_events": 5, "highest_risk": 73, "confirmed": 3, "watchlist": 1, "conflicting": 1 },
  "events": [
    { "id": "EVT-001", "name": "智利锂矿罢工", "type": "production_disruption",
      "severity": 4, "confidence": "confirmed", "risk_score": 73 }
  ],
  "map": {
    "points": [ { "lat": -23.65, "lng": -70.40, "label": "智利锂矿", "score": 73 } ],
    "arcs": [ { "from": "智利", "to": "中国", "material": "碳酸锂" } ]
  },
  "markets": {
    "lithium_etf": { "value": 45.20, "change": "+3.2%", "status": "available" },
    "spot_price": { "status": "unavailable", "note": "数据源连接超时" }
  },
  "report": {
    "status": "llm_generated",
    "content": "# 产业链风险研判报告\n## 一、事件概述\n..."
  },
  "data_quality": {
    "gaps": [ { "field": "锂现货价格", "reason": "数据源连接超时" } ],
    "inferred_nodes": [ { "event": "智利锂矿罢工", "node": "智利", "method": "精确匹配" } ]
  }
}
\end{verbatim}

从示例中可以看出，当市场数据（spot\_price）不可用时，接口不会隐藏这个信息，而是明确返回 status: "unavailable" 及原因。Dashboard 的前端组件接收到这个状态后，会在对应面板展示"数据暂不可用"的提示，而不是显示一个空白面板或过时数据。这种设计背后的理念是：在风险研判场景中，"知道数据不可用"本身就是一种有价值的信息——它告诉用户这个判断维度目前没有被纳入考虑。如果把缺失数据默默地忽略掉，用户可能误以为系统已经评估了所有维度。

\section{工程实现状态}

当前系统已完成以下工程内容：

\textbf{数据层面：} 
\begin{itemize}
    \item 锂电池产业链种子知识图谱：覆盖 5 类节点（资源国 8 个、材料 15 种、产业环节 6 个、企业 10+、物流通道 3 个）和 4 类关系边（supplies、depends\_on、processes、transports），共约 120 个节点和 200 条关系边。采用 JSON 文件存储，加载后以邻接表形式保存在内存中，支持图谱路径搜索。
    \item 来源信誉注册表：约 40 个信息来源的分类和信誉度赋值，以 JSON 配置文件维护。
    \item 地理位置映射表：约 60 条国家/地区/港口到经纬度的映射，用于 Dashboard 地图点位。
    \item 演示案例数据：30--50 条覆盖政策、罢工、航运、价格等场景的预置事件数据，支持离线模式运行。
    \item 多源 RSS 采集管线：基于 Node.js 的定时抓取模块，从约 15 个行业新闻源、政策公告源和物流数据源采集原始文本。
\end{itemize}

\textbf{算法层面：} 
\begin{itemize}
    \item 事件抽取管线：规则模块（正则 + 词典）用 JavaScript 实现，约 800 行代码。LLM 补全省调用 API 完成。Schema 校验器用 JSON Schema 定义字段约束。
    \item 四级实体链接器：别名表约 150 条映射，以 Map 数据结构驻留内存，查询时间复杂度 O(1)。
    \item 置信度计算与状态机：纯函数实现，输入事件对象输出置信度分数和状态标签，不依赖外部服务。
    \item 图谱传播路径搜索：BFS 实现，最大深度 5 跳，路径排序使用加权评分函数。平均单次搜索时间约 5--15ms（当前图谱规模下）。
    \item 可解释评分模型：5 因子线性加权，每个因子的计算函数独立实现，便于后续单独替换或调参。
    \item 证据链 Prompt 组装器：模板引擎 + 条件判断，根据事件数量、数据可用性等动态生成结构化 Prompt。
\end{itemize}

\textbf{展示层面：} 基于 HTML/CSS/JS 的 Dashboard，包含 2D/3D 地图组件（基于 Leaflet + Globe.gl）、事件列表面板、市场指标面板（K 线图和折线图）、AI 报告组件（Markdown 渲染）、数据缺口提示组件。前端通过 REST API 与后端通信。

\textbf{演示保障层面：} 系统设计了三级降级机制。第一级：有网络时调用 LLM API 生成完整报告。第二级：LLM 调用失败时自动切换为 fallback 模板报告。第三级：整个后端不可用时，Dashboard 加载预置的缓存数据文件，前端仍然可以展示完整的 Demo。此外，Dashboard 的每个数据接口字段都有默认值——地图没有最新点位就展示缓存点位，市场数据取不到就显示"数据暂不可用"，不会因为某个模块异常导致页面空白。

\section{验证方式}

初赛阶段的验证采用案例驱动方式。我们准备了 4--6 个典型风险场景作为验证案例：

\begin{itemize}
    \item 印尼镍矿出口限制政策（政策变化类）
    \item 智利锂矿工人罢工（生产中断类）
    \item 红海航运安全事件（物流中断类）
    \item 刚果金钴矿政策调整（资源国风险类）
    \item 欧盟电池法案更新（法规变化类）
\end{itemize}

每个案例从 7 个环节进行检查：事件抽取字段完整率、实体映射准确性、置信度状态合理性、传播路径是否符合产业链常识、风险评分与人工判断的一致性、Dashboard 展示与后端数据的一致性、报告中证据引用的可追溯性。

\begin{longtable}{p{2.4cm}p{3.5cm}p{7.0cm}}
\caption{验证矩阵}\label{tab:evaluation}\\
\toprule
验证对象 & 验证指标 & 验证方式 \\
\midrule
\endfirsthead
\toprule
验证对象 & 验证指标 & 验证方式 \\
\midrule
\endhead
事件抽取 & 字段完整率 & 检查事件类型、时间、地点、材料、证据是否完整。 \\
实体映射 & 节点匹配准确率 & 人工核对事件是否落到正确的图谱节点上。 \\
置信度状态 & 状态一致性 & 检查 confirmed / watchlist / conflicting 的判定是否合理。 \\
传播路径 & 路径合理性 & 逐一判断路径是否符合产业链真实传导关系。 \\
风险评分 & 等级一致性 & 与人工判定的高中低风险进行对照。 \\
报告内容 & 证据可追溯性 & 报告中每条结论都能对应到具体事件、路径或评分字段。 \\
Dashboard & 数据一致性 & 检查地图点位、事件列表、弧线是否与后端输出一致。 \\
\bottomrule
\end{longtable}

\section{与现有技术方案的差异}

我们做了一个对比来明确自己的定位：

\begin{longtable}{p{3.0cm}p{3.2cm}p{3.6cm}p{3.8cm}}
\caption{技术方案对比}\label{tab:comparison}\\
\toprule
方案类型 & 能做什么 & 做不到什么 & 我们的改进 \\
\midrule
\endfirsthead
\toprule
方案类型 & 能做什么 & 做不到什么 & 我们的改进 \\
\midrule
\endhead
普通新闻聚合 & 聚合标题和摘要 & 无法判断信息真假和产业链影响范围 & 事件结构化 + 图谱传播推理 \\
通用大模型问答 & 回答开放域问题 & 容易把推测当成事实，无法给出可解释评分 & 证据链约束 + 结构化评分 \\
普通 RAG & 检索相关文本片段 & 无法表达事件状态变更和跨文本结构关系 & 时序状态机 + 产业链图谱 \\
传统舆情系统 & 关键词监测、热度曲线 & 缺少行业结构感知和风险量化能力 & 可解释评分 + 行业图谱 \\
单一可视化看板 & 展示已有指标 & 不具备推理能力和证据追溯机制 & 全链路推理 + 数据缺口显式展示 \\
\bottomrule
\end{longtable}

核心差异可以概括为一句话：\textbf{我们的系统不是在"找相关信息"，而是在"做产业链风险推理"。} 前者是检索问题，后者是结构推理问题。

这五种方案的对比维度可以从三个层面来看：

\textbf{第一层：信息处理方式。} 新闻聚合和舆情系统处理的是"关键词"和"热度"——它们能告诉你"有多少人在讨论锂矿"，但回答不了"这条消息是真的还是假的""它会影响哪些产业链环节"。RAG 系统处理的是"文本片段"——它能找到相关段落，但段落本身不包含产业链结构。我们的系统处理的是"结构化事件 + 图谱关系"——每条事件被转化为可计算的字段，每个实体被固定到图谱节点上，推理在结构层面进行而非文本层面。

\textbf{第二层：推理能力。} 通用大模型具有强大的语言理解和生成能力，但它有两个根本限制：没有可靠的内部事实核查机制（它无法区分"这件事被官方确认了"和"这件事有人在推特上说了"），也没有产业链结构知识（它不知道正极材料在电芯的上游）。我们的系统把这两个限制通过架构设计来解决——可信度评估模块负责事实核查，知识图谱负责结构知识，大模型只做语言组织。

\textbf{第三层：可解释性。} 大模型给出的风险判断本质上是黑箱——你无法知道它为什么给某个事件打了高风险。舆情系统的热度曲线可以解释（因为是统计量），但它解释不了产业链影响。我们的风险评分是显式可拆解的——每个分数都能追溯到五个因子的具体贡献，每条传播路径都能在图谱上画出来，每个置信度状态都能关联到具体的来源和证据。

\section{当前局限与后续计划}

坦诚地说，当前系统仍处在原型阶段，以下方面需要持续完善：

\begin{itemize}
    \item \textbf{图谱覆盖有限。} 当前的种子图谱覆盖了核心矿产（锂、镍、钴、石墨等 8 种）、材料（碳酸锂、氢氧化锂、硫酸镍等 15 种）、产业环节（正极材料、电芯、电池包、整车、储能 6 个）和部分物流通道（红海、马六甲、苏伊士 3 个），但尚未扩展到具体矿山、供应商和企业级的贸易关系。这意味着系统目前能回答"锂矿罢工会影响碳酸锂和正极材料"，但回答不了"这个罢工具体影响哪家企业的哪个工厂"——后者需要企业级的供应链数据，这类数据通常不公开，需要通过商业数据源或企业合作获取。初赛阶段支撑案例演示够用，后续如果进入产品化阶段，需要引入图数据库（如 Neo4j）来管理更大规模（预计 10 万+ 节点）的产业链关系。
    
    \item \textbf{评分权重待校准。} 五个因子的权重（0.35/0.25/0.20/0.15/0.05）当前基于经验设定。我们没有理由认为这组权重在任何场景下都是最优的——不同行业、不同类型的风险，因子的重要性可能不同。比如对于金融投资者，时间新鲜度可能更重要（市场对新闻的反应速度很快）；对于供应链管理者，节点关键性可能更重要（关键节点断供的影响远大于非关键节点）。后续计划通过构建 200--300 条人工标注的验证集（每条标注真实风险等级），然后用网格搜索 + 交叉验证的方式找出最优权重组合，并针对不同风险场景训练不同的权重配置。
    
    \item \textbf{事件抽取依赖 LLM。} 尽管有规则兜底（候选实体识别和 schema 校验都是规则实现），但字段补全和复杂文本理解仍依赖大模型 API 调用。这带来几个实际问题：API 调用有延迟（通常 1--3 秒，高峰期可能更长）、有成本（按 token 计费）、有网络依赖（断网时无法调用）。如果批量处理大量新闻，LLM 调用的延迟会成为瓶颈。后续的改进方向是：积累 500--1000 条经过 LLM 补全的事件样本，用这些数据微调一个小参数量模型（如 7B 级别的开源模型），替代通用大模型的 API 调用。微调模型可以本地部署，延迟降至毫秒级，且没有 API 调用成本。
    
    \item \textbf{市场指标需要更稳定的数据源。} 当前接入的材料价格和航运数据通过公开或低成本的 API 获取，稳定性不够理想——部分免费 API 有调用次数限制，部分数据源在国内访问不稳定。这导致市场指标面板在某些演示中显示"数据不可用"，影响展示效果。后续需要接入商业级数据源（如 S\&P Global、Bloomberg 的 API），或者与行业研究机构合作获取更稳定的数据。
\end{itemize}

\section{结论}

这个系统不是"让大模型读新闻然后给结论"——那是很多人第一时间能想到的做法，但我们认为那条路走不通，理由有三：

第一，大模型没有可靠的机制区分"确认的事实"和"网络传言"。它会把权威公告和推特帖子一视同仁地当作输入文本处理，无法判断哪个更可信。第二，大模型没有产业链结构知识。它可能知道"锂是电池的重要原料"，但它不知道从锂矿到电池包之间具体有哪些加工环节、每个环节的全球产能分布、以及哪些节点是关键瓶颈。第三，大模型的风险判断不可拆解。它给一个事件打"高风险"时，你无法知道它是基于什么理由——是事件本身严重，还是受影响环节关键，还是因为信息来源可靠？评分不可拆解就意味着不可复核、不可调试。

我们的做法绕开了这三个问题：

用事件结构化把非结构化信息变成可计算对象——这样后续模块不需要"理解"自然语言，只需要读取字段值。用产业链图谱建模传导关系——这样系统能回答"影响哪些环节"而不需要依赖文本中提到。用置信度评估过滤噪声——这样系统能区分"政府公告"和"推特传言"。用图搜索找出传播路径——这样系统能计算"锂矿罢工"到"电池包"之间的具体传导链。用加权评分给出可解释的风险量化——这样每个分数都能拆解出五个因子的具体贡献。最后让大模型在证据链约束下生成报告——这样大模型只负责语言组织，不负责风险判断。

每一步都是在\textbf{给大模型划边界}，而不是\textbf{指望大模型做判断}。我们认为，对于一个需要落地到真实产业链风险场景的系统来说，这种设计思路比"端到端大模型"更可靠——因为它每个环节都可以单独验证、单独调试、单独改进，而不是把全部赌注押在一个黑箱模型上。

\end{document}
